\subsection{Use Case UC-11}\label{use-case-uc-11}

\begin{longtable}[]{@{}|>{\raggedright\arraybackslash}p{0.27\columnwidth}|>{\raggedright\arraybackslash}p{0.67\columnwidth}|@{}}
\hline
\rowcolor{TableBlue}
\tableheader{Use case ID} & \tableheader{UC-11} \\ \hline
\emph{Use case name} & Quản lý danh mục vật phẩm \\ \hline
\emph{Brief description} & Admin thêm, sửa hoặc ẩn các danh mục vật phẩm (ví, thẻ, điện tử…). \\ \hline
\emph{Primary actor} & Admin \\ \hline
\emph{Secondary actors} & Không có \\ \hline
\emph{Priority} & Must \\ \hline
\emph{Related requirements} & FR-08 \\ \hline
\emph{Trigger} & Admin mở mục “Quản lý danh mục”. \\ \hline
\emph{Pre-conditions} & Admin đã đăng nhập. \\ \hline
\emph{Post-conditions} & Danh sách danh mục được cập nhật trong form đăng bài và bộ lọc tìm kiếm. \\ \hline
\emph{Main flow} & 1. Admin mở “Quản lý danh mục”. 2. Hệ thống hiển thị các danh mục kèm số bài trong mỗi danh mục. 3. Admin nhấn “Thêm” và nhập tên danh mục. 4. Hệ thống kiểm tra tên. 5. Hệ thống lưu danh mục mới; danh mục này xuất hiện trong form đăng bài và bộ lọc. \\ \hline
\emph{Alternative flows} & A1 (tại bước 3): Admin chọn “Sửa” một danh mục và nhập tên mới; tiếp tục từ bước 4. A2 (tại bước 3): Admin chọn “Ẩn” một danh mục. Hệ thống ẩn danh mục khỏi form đăng bài, nhưng các bài cũ vẫn giữ danh mục đó. \\ \hline
\emph{Exception flows} & E1 (tại bước 4): Tên để trống hoặc trùng với danh mục đã có. Hệ thống báo lỗi. \\ \hline
\emph{Business rules} & Không có \\ \hline
\emph{Non-functional constraints} & Không có \\ \hline
\end{longtable}
